\newcommand{\slidemode}{1}
% Modified by LianTze Lim to work with fontspec/xelatex
    % Choose for notes or slides
    \if \slidemode 1
        \documentclass{beamer}
        \else
        \documentclass{article}
        \usepackage{beamerarticle}
        \usepackage[a4paper,margin=2cm]{geometry}
        \usepackage{multicol}
    \fi 
\usepackage{amsfonts,amsmath,amssymb,amsthm,amscd}
\usepackage[style=authoryear]{biblatex}
\renewcommand*{\nameyeardelim}{\addcomma\addspace}
\AtBeginBibliography{\tiny}

\usetheme{DarkConsole}
\setbeamertemplate{theorems}[normal font]

% Packages
\usepackage{hyperref}
\usepackage{graphicx}
\usepackage{multimedia} % provides \movie
\usepackage{tikz,pgfplots}
\usetikzlibrary{arrows.meta,calc}
\pgfplotsset{compat = 1.18}
\usepgfplotslibrary{fillbetween}
\usepackage{physics}
\usepackage{siunitx}
\usepackage{polynom}
\usepackage{biblatex}
\usepackage[outline]{contour} % glow around text
\usetikzlibrary{angles,quotes,intersections,fillbetween} % for pic
\newcommand{\myscale}{0.7}
\contourlength{1.3pt}

\tikzset{>=latex} % for LaTeX arrow head
\usepackage{xcolor}
\tikzstyle{vector}=[->,very thick,xcol]
\tikzstyle{mydashed}=[dash pattern=on 2pt off 2pt]
\def\tick#1#2{\draw[thick] (#1) ++ (#2:0.1) --++ (#2-180:0.2)} 

% helper: open-dot macro (hollow circle)
  \newcommand{\opendot}[2]{%
    \filldraw[fill=white,draw=gray,thick] (#1,0) circle[radius=2.2pt];
    \node[below=4pt] at (#1,0) {#2};
  }
  \newcommand{\closedot}[3][magenta]{%
    \filldraw[fill=#1,draw=#1,thick] (#2,0) circle[radius=2.2pt];
    \node[below=4pt] at (#2,0) {#3};
  }
\usepackage[brazilian]{babel}
\usepackage{diagbox}
\usepackage{booktabs}

% collor-blind friendly pallete
\definecolor{figcolor1}{RGB}{0,70,70}
\definecolor{figcolor2}{RGB}{116,233,233}
\definecolor{figcolor3}{RGB}{146,0,0}
\definecolor{color1}{RGB}{0,70,70}
\definecolor{color2}{RGB}{116,233,233}
\definecolor{color3}{RGB}{146,0,0}


% A counter, since TikZ is not clever enough (yet) to handle
% arbitrary angle systems.
\newcount\mycount
\usepackage[percent]{overpic}
\usepackage{xcolor}
\usepackage{adjustbox}
\usepackage{bbold}
\newcommand{\rest}{\upharpoonright}

% symbols
\newcommand{\F}{\mathcal{F}}
\newcommand{\C}{\mathcal{C}}
\newcommand{\Open}{\mathcal{O}}
\newcommand{\A}{\mathcal{A}}
\newcommand{\B}{\mathcal{B}}
\newcommand{\G}{\mathcal{G}}
\newcommand{\M}{\mathcal{M}}
\newcommand{\N}{\mathbb{N}}
\newcommand{\Q}{\mathbb{Q}}
\newcommand{\meta}[1]{\overset{\text{meta}}{#1}}
\newcommand{\val}[1]{\llbracket #1 \rrbracket}
\newcommand{\valin}[3]{
	\sup_{#3 \in dom(#2)}\val{#3 = #1}#2(#3)}
\newcommand{\valsub}[3]{
	\inf_{#3 \in dom(#1)}(#1(#3)\implies \val{#3 \in #2})}
\newcommand{\valeq}[2]{\val{#1 \subset #2}\val{#2 \subset #1}}
% epimorphism arrow
\newcommand{\rightarrowdbl}{\rightarrow\mathrel{\mkern-14mu}\rightarrow}
\newcommand{\epimorphism}[2][]{%
	\xrightarrow[#1]{#2}\mathrel{\mkern-14mu}\rightarrow
}
\newcommand{\catname}[1]{\textnormal{\textsc{#1}}}

\newcommand{\R}{\mathbb{R}}
\newcommand{\D}{\mathcal{D}}
\newcommand{\Part}{\mathbb{P}}
\newcommand{\sequence}[1]{\langle #1 \rangle}
\DeclareMathOperator{\inter}{int}
\newcommand{\interior}[2]{\inter_{#2} (#1)}
\DeclareMathOperator{\clos}{cl}
\DeclareMathOperator{\diam}{diam}
\newcommand{\closure}[2]{\clos_{#2} (#1)}
\newcommand{\st}{\,:\,}
\newcommand{\dem}{\emph{dem.:}\\}
\renewcommand{\d}{\mathrm{d}}

% Theorem environments
\numberwithin{equation}{section}
% Theorem environments
\newtheorem{teor}{Teorema}[section]
\newtheorem{lema}{Lema}[section]
\newtheorem{prop}{Proposição}[section]
\newtheorem{properties}{Propriedades}[section]
\newtheorem{cor}{Corolário}[section]  
\newtheorem{defin}{Definição}[section]
\newtheorem{ex}{Exemplo}[section]
\newtheorem*{theorem*}{Teorema}
\newtheorem*{proposition*}{Proposição}
\newtheorem*{corollary*}{Corolário}
\newtheorem*{lemma*}{Lema}
\theoremstyle{remark}
\newtheorem{rmk}{Observação}[section]
\newtheorem*{question}{Problema}
\newtheorem{exerc}{Exercício}

\title{\texttt{Algebrizando a lógica}}
\author{Violeta Mar}
\date{}
\addbibresource{bibliography.bib}

\begin{document}

\begin{frame}{}
\maketitle
\end{frame}

\begin{frame}{}
\frametitle{O que é lógica?}
\pause
\begin{itemize}
    \item Lógica (como uma ciência, uma área do conhecimento) é o estudo de um tipo muito específico de pensamento.\pause
    \item Pensamento racional? Analítico? Dedutivo? Formal? Mecânico? \pause
    \item A lógica matemática 'captura' este tipo de pensamento em definições matemáticas (sistemas axiomáticos e dedutivos, linguagens formais, etc). \pause
    \item O cálculo proposicional e a lógica de primeira ordem são os sistemas lógico-matemáticos mais comuns e difundidos, mas também temos muitos outros (lógicas não clássicas, lógicas modais, lógicas paraconsistentes, etc). \pause
    \item Por essas coisas serem objetos matemáticos, podemos simplesmente fazer matemática sobre elas!
\end{itemize}
\end{frame}

\begin{frame}{}
\frametitle{O Cálculo Proposicional}
\pause
\begin{itemize}
\item O cálculo proposicional é uma linguagem formal cujo alfabeto é composto de enumeráveis símbolos $p_1,p_2, p_3, \dots$ e os símbolos conectores $\wedge, \vee, \lnot, \rightarrow$ \pause
\item As 'palavras' formadas por este alfabeto de forma 'correta' são chamadas de fórmulas, e são coisas do tipo $p_2 \wedge \lnot p_1 \rightarrow p_3 \vee p_5$ \pause
\item Usar como uma linguagem já é útil pra matemática! \pause
\item Mas fazemos mais: definimos algumas fórmulas como especiais (por exemplo, as fórmulas $p \rightarrow p$, $p \wedge \lnot p$ e $\lnot(p \vee \lnot p)$ parecem representar 'verdades'), que vão ser as chamadas 'tautologias' \pause
\item Também teremos um sistema dedutivo onde definimos quando que um conjunto $T$ de fórmulas deduz uma outra fórmula $\alpha$. (por exemplo, a partir das fórmulas $p \rightarrow q$ e $p$, podemos deduzir $q$) \pause
\item De novo, estas são práticas comuns no fazer matemático. 
\end{itemize}
\end{frame}

\begin{frame}{}
\frametitle{A Lógica de Primeira Ordem}
\pause
\begin{itemize}
\item A lógica de primeira ordem é uma extensão do cálculo proposicional onde agora temos também símbolos que irão representar variáveis, constantes, funções, relações e quantificadores. \pause
\item As fórmulas vão ser coisa do tipo $\exists x ((f(x,c) = d) \vee (R(x) \rightarrow x = d))$ \pause
\item Agora nossa linguagem consegue falar de basicamente qualquer tópico matemático. \pause
\item Também vamos ter aqui as fórmulas especiais (as que são sempre verdade), e um sistema dedutivo (que é basicamente Modus Ponens também). 
\end{itemize}
\end{frame}

\begin{frame}{}
\frametitle{Semântica e resultados}
\end{frame}

\begin{frame}{}
\frametitle{O que é álgebra?}
\begin{itemize}
    \item Álgebra é o estudo matemático de conjuntos com operações. \pause
    \item Operações: somar um número com outro número, compor uma função com outra função, aplicar uma simetria após a outra... \pause 
    \item Estudamos as 'propriedades gerais' que um conjunto com uma (ou mais que uma) operação deve satisfazer (geralmente, sempre supondo que ele satisfaça certas propriedades, colocadas em equações).\pause
    \item Botar exemplo
\end{itemize}
\end{frame}

\begin{frame}{}
\frametitle{A álgebra da lógica}
Mas em lógica temos operações no conjunto de fórmulas: $\vee, \wedge, \exists$... Qual é a estrutura algébrica aqui?
\end{frame}

\begin{frame}{}
\frametitle{Álgebra de Lind Tarski}
\end{frame}

\begin{frame}{}
\frametitle{Álgebra de Boole}
\end{frame}

\begin{frame}{}
\frametitle{Álgebra de Halmos}
\end{frame}

\begin{frame}{}
\frametitle{Traduzindo!}
\end{frame}

%lógicas não clássicas?

\end{document}
